\section{Related Work}
\label{sec:related_work}

\para{LLM trading agent architectures.}
The design of LLM trading agents has progressed rapidly from simple prompt-based systems to sophisticated multi-agent architectures. \finmem~\citep{yu2023finmem} introduced a layered memory architecture with shallow (daily), intermediate (quarterly), and deep (annual) memory, achieving strong short-term results on individual stocks. \fincon~\citep{yu2024fincon} extended this with a manager-analyst hierarchy and Conceptual Verbal Reinforcement, showing particular strength in bearish markets. \tradingagents~\citep{xiao2025tradingagents} mirrors the organizational structure of real trading firms with specialized analyst, researcher, and risk management agents. Despite these architectural advances, all three systems---and every other published LLM trading agent we surveyed---operate exclusively at daily decision frequency.

\para{Rigorous evaluation of LLM trading.}
Recent work has challenged the optimistic results reported in earlier studies. \finsaber~\citep{li2026finsaber} introduced a bias-mitigated backtesting framework that evaluates LLM strategies over 20 years across 100+ stocks, finding that strategies appearing profitable in short-term evaluations (6 months, 4--8 stocks) consistently underperform Buy-and-Hold at scale. The key diagnosed failure mode is \emph{regime miscalibration}: LLMs are overly conservative in bull markets and overly aggressive in bear markets. \deepfund~\citep{deepfund2025} complemented this with live (non-backtested) evaluation, finding that only 1 of 9 LLMs achieved positive returns over 24 trading days---and that model succeeded via conservative cash management rather than superior prediction. \stockbench~\citep{chen2025stockbench} constructed a contamination-free benchmark on post-training-cutoff data, finding that most LLMs only marginally beat Buy-and-Hold while providing better risk management. These findings collectively suggest that daily LLM trading faces fundamental challenges, but none of these studies considered alternative decision frequencies.

\para{Reinforcement learning approaches.}
An alternative line of work uses reinforcement learning to fine-tune LLMs for trading. \flagtrader~\citep{li2025flagtrader} showed that a 135M-parameter LLM fine-tuned with PPO outperforms GPT-4 in agentic frameworks, using Sharpe ratio change as the RL reward. While RL-based approaches address some limitations of pure prompting, they still operate at daily frequency and require substantial computational resources for training. Our work is complementary: we investigate whether the \emph{decision horizon} itself is a key variable, independent of the agent architecture.

\para{Multi-frequency analysis in traditional finance.}
The idea that trading performance depends on decision frequency is well-established in traditional finance. Momentum strategies show frequency-dependent returns~\citep{jegadeesh1993returns}, and optimal rebalancing frequency varies by asset class and market regime. However, this dimension has not been explored for LLM-based agents. Our work bridges this gap by applying the lens of frequency analysis to the emerging field of LLM trading.

\para{Benchmarks and surveys.}
\investorbench~\citep{li2024investorbench} provides a benchmark testing 13 LLMs across stocks, crypto, and ETFs, finding that ETF trading---the most strategic, long-term task---was hardest for most LLMs but improved most with capable models. \citet{ding2024survey} surveyed 14 LLM trading systems and found a median backtesting period of only 1.3 years, with trading costs almost universally ignored. Both observations motivate our investigation of longer-horizon strategies that naturally reduce trading frequency and costs.
